\section{BIG-Bench Hard: del promedio a la heterogeneidad por tarea}

Una sola demo de arithmetic puede ser simplemente cherry-pick, por eso las diapositivas pasan a BIG-Bench Hard (BBH): de BIG-Bench se seleccionan las tareas difíciles en las que los modelos de entonces quedaban por debajo del average human rater, y luego se compara de forma sistemática answer-only contra CoT. Este diseño descompone la pregunta «¿puede hacer razonamiento complejo?» en varias task mechanics, como navegación, ordenamiento, lógica y seguimiento de estado.

\subsection{Cómo define el benchmark lo «hard»}

La dificultad de BBH proviene de las condiciones de construcción, no de una propiedad permanente: en las tareas seleccionadas, ningún modelo superó al average human en los resultados originales de BIG-Bench. El prompt incluye task description, worked examples, options y reasoning trace; por lo tanto, lo que se evalúa es si el modelo puede inferir el algorithm a partir de la descripción en lenguaje natural y de los exemplars.

\lecturefigure{slide-24.jpg}{Selección de tareas de BIG-Bench Hard y estructura del prompt con CoT}{Diapositiva 24 del material oficial del curso; Suzgun et al., 2022}

\begin{knowledgebox}{Lectura de la figura: dos ejemplos que miden estados intermedios distintos}
Navigation exige actualizar de forma continua la dirección y las coordenadas, y al final decidir si se volvió al origen; word sorting exige comparar las letras iniciales, manejar prefijos idénticos y mantener el orden de salida. Ninguna de las dos tareas es misteriosa para una persona, pero ambas requieren que el modelo mantenga un estado coherente a lo largo de varios tokens. El valor de CoT está en escribir explícitamente las actualizaciones de estado, en lugar de adivinar solo al final un yes/no o una lista.
\end{knowledgebox}

\subsection{Además del promedio, hay que ver cuántas tareas cruzan la línea humana}

Si solo se reporta la aggregate accuracy, unas pocas tareas con grandes ganancias pueden ocultar muchos fracasos. Las diapositivas reportan a la vez la performance promedio y la «proporción de tareas que superan al average human rater», y muestran la heterogeneidad con un mapa de tareas en rojo y azul. Las dos metric responden, respectivamente, a «cuánto mejora el total» y «cuánto se amplía la cobertura de capacidades».

\lecturefigure{slide-25.jpg}{Panorama de resultados de BBH: rendimiento promedio, número de tareas que superan al humano y mapa por tarea}{Diapositiva 25 del material oficial del curso; Suzgun et al., 2022}

\begin{knowledgebox}{Lectura de la figura: primero entender el subset bias, luego la ganancia}
BBH selecciona de origen tareas en las que los modelos anteriores quedaban por debajo del humano, así que el baseline answer-only bajo es by construction. Al agregar CoT, tanto el promedio de davinci-002 como el número de tareas que superan al average-human aumentan; el mapa rojo y azul de abajo muestra que la ganancia no es uniforme. La conclusión debe ser que CoT desbloquea la mayoría de las tareas difíciles, pero no todas, y no que «el modelo supera al humano en todo».
\end{knowledgebox}

\teachervoice{En clase se explican una por una las dos metric y se señala que el average human ronda 67, mientras que el best human es más alto. El ponente no presenta «la mayoría de las tareas supera al average human» como un nivel humano general, sino que lo restringe al conjunto de problemas, el prompt y el método de calificación de BBH.}

\subsection{Scaling: el delta positivo de CoT tiene su propio umbral}

La siguiente figura descompone los resultados del modelo más grande en varias escalas. Si CoT solo añadiera información al prompt, todos los modelos deberían obtener una ganancia relativa similar; sin embargo, las curvas muestran que los modelos pequeños no necesariamente obtienen un delta positivo, y solo alrededor de davinci-002 / PaLM 62B se observa un beneficio estable en el aggregate. Este scale sweep es fundamental desde el punto de vista metodológico, porque convierte un «caso de éxito en el modelo más grande» en una prueba de la interacción entre la intervention y la base capability.

\lecturefigure{slide-26.jpg}{Scaling de CoT en BBH: la ganancia positiva requiere un tamaño de modelo suficiente}{Diapositiva 26 del material oficial del curso; Suzgun et al., 2022}

\begin{knowledgebox}{Lectura de la figura: el delta identifica mejor que la puntuación absoluta el umbral de un método}
En cada panel se compara no-CoT con CoT para el mismo modelo, en lugar de comparar solo las puntuaciones absolutas de modelos distintos. En los modelos pequeños, la reasoning trace puede introducir errores adicionales; solo los modelos grandes aprovechan la trace para descomponer la tarea, y entonces aparece una diferencia positiva estable. El umbral de la figura es una posición empírica para una model family y un task aggregate específicos; no debe escribirse como una regla fija de 62B para todo sistema con CoT.
\end{knowledgebox}

\subsection{Emergence: No-CoT sigue plano y la curva de CoT empieza a subir}

El relato más contundente de la clase surge del contraste: la performance answer-only sigue casi flat en las escalas evaluadas, mientras que CoT hace subir de forma notable a los modelos más grandes. Aquí «desbloquear» significa que la interface permite que una capacidad composicional ya existente, pero no invocable directamente, se refleje en la metric; no demuestra que, tras el preentrenamiento, el modelo carezca por completo del conocimiento relevante.

\lecturefigure{slide-27.jpg}{Emergence de CoT en BBH: answer-only se mantiene plano y CoT abre el crecimiento del rendimiento}{Diapositiva 27 del material oficial del curso; Suzgun et al., 2022}

\begin{knowledgebox}{Lectura de la figura: el prompt puede cambiar la frontera de capacidad que observamos}
Si solo se evalúa no-CoT, los investigadores podrían concluir que «seguir escalando no produce avances»; al agregar reasoning exemplars, la misma familia de modelos muestra ganancias asociadas a la escala. Esto indica que la capability evaluation no puede separarse del elicitation method. La figura respalda que «CoT eleva la capacidad medible en este escenario», pero no demuestra que toda flat curve pueda desbloquearse con un mejor prompt.
\end{knowledgebox}

\begin{importantbox}{Al evaluar la capacidad de un modelo hay que reportar también el elicitation budget}
El modelo, el prompt, los few-shot examples, el generated-token budget, la sampling strategy y las external tools definen en conjunto una prueba de capacidad. Reportar solo el parameter count y la final accuracy lleva a registrar una interface failure como capability absence, y también oculta una inference-time search costosa bajo la idea de que «el modelo en sí es más inteligente».
\end{importantbox}

\subsection{Resumen de la sección}

BBH extiende CoT de un ejemplo aislado a un conjunto de tareas difíciles. El promedio, la proporción de tareas que superan al humano y el mapa por tarea muestran en conjunto ganancias heterogéneas; el scale sweep indica además que CoT es en sí una emergent prompting technique. Por ello, las curvas de capacidad deben interpretarse junto con el elicitation method.

